\section{实践建议与系统运维}
\subsection{Autonomy 监控仪表盘}
讲者建议建立包含 reward accuracy、reset success、goal coverage 三个指标的仪表盘，用于日常 review。
\begin{tabularx}{\textwidth}{lX}
\toprule
维度 & 指标 \\
\midrule
Reward fidelity & classifier precision/recall, clip similarity drift； \\
Reset coverage & 近期 forward-backward > 80\% 成功率； \\
Goal expansion & 活跃 goal 数与每周新增 goal； \\
\bottomrule
\end{tabularx}

\subsection{跨阶段反馈循环}
每运行一次 Autonomy experiment，记录：
\begin{enumerate}
\item 当前 reward 模型与 safety penalty 版本；
\item Reset policy 的 failure pattern；
\item 新目标/技能的添加、旧目标的 retire 决策；
\end{enumerate}
这些信息形成一个 cross-team log，便于 future replication。

\begin{knowledgebox}{Feedback Loop 模板}
1. Pre-scale：确认 reward/safety/checkpoint 状态；\\
2. Scale：运行 policy，同步 metrics；\\
3. Post-scale：复盘 metrics drift，如需 rollback 立即触发。
\end{knowledgebox}

\subsection{本章小结}
实践层面，Autonomy 需要仪表盘、跨阶段反馈以及回退计划，确保研究结果可反复部署。
\newpage

